 % !TeX root = book.tex
\chapter{The Trinity as a web-native symbolic paradigm}\label{sec:paradigm}

The preceding chapters have introduced the elements of the Prolog Trinity ecosystem. We began with Web Prolog as a carefully delimited profile of Prolog, extended with Erlang-style message passing and concurrency-oriented semantics. We then examined Prolog agents -- actors and nodes -- as concrete computational entities, each equipped with clearly defined capabilities and lifecycle properties. From there, we explored the Prolog Web itself: the network that arises when such entities are allowed to interact across transport boundaries, using both sequential query-response patterns and long-lived concurrent conversations. Subsequent chapters introduced reusable behavioral patterns, statechart actors, and the positioning of the Trinity in relation to the Semantic Web and contemporary AI architectures.

Taken individually, these components may appear as incremental extensions of existing technologies. Web Prolog resembles Prolog with additional concurrency primitives. The actor abstraction draws inspiration from Erlang. The sequential layer echoes long-established patterns of remote querying. Statecharts have been known for decades. 

Yet when considered together, these elements form a coherent architectural stance -- one that may reasonably be viewed as a candidate paradigm for executable symbolic systems on the Web. And perhaps not only symbolic. Neuro-symbolic integration is a growing area across the AI community. The Web is becoming agentic, and it seems evident that a language that can serve both as a logic programming language and as an agent programming language may have something to contribute.

Before developing the argument, one point of orientation. The Trinity is not proposed as a competing system but as scaffolding -- an ecosystem skeleton -- around which existing Prolog systems could collaborate without excluding anyone and without making any system less powerful. The community and adoption implications of this stance are developed in Chapter~\ref{sec:whats-in-it-for-the-prolog-community}.

With that framing in place, this chapter asks a disciplined question:

\begin{quote}
\normalsize What conceptual pattern emerges when logic programming, actor-based concurrency, Web transport semantics, and structured behavioral abstractions are deliberately integrated into a single design?
\end{quote}

\noindent The answer, we will argue, is a paradigm characterised by a small number of interlocking commitments: a foundational separation between sequential and concurrent computation; a semantic rather than merely syntactic notion of portability; the agent as the primary unit of structure; a logical foundation that survives distribution; and the treatment of interaction governance as a first-class design dimension.

These commitments are mutually constraining. The language profile determines what can be exposed safely and portably; the agent model determines which long-lived behaviours can be made robust; the network substrate determines which forms of composition are realistic at Web scale. It is precisely this mutual constraint that prevents the proposal from collapsing into ``Prolog over HTTP'', ``actors in Prolog'', or ``a distributed Prolog database''. No single strand can be redesigned in isolation without consequences for the others.


%%% ───────────────────────────────────────────────────────────────
\section{The two partitions as a foundational principle}\label{sec:two-partitions}

One of the central structural decisions in the Trinity architecture is
the separation between two distinct partitions of computation: the
sequential and the concurrent.

The \emph{sequential partition} is the world of queries and relations. A
client poses a query, a node evaluates it against a clause database, and
solutions are produced nondeterministically. Over the Web, this
partition manifests as remote querying, paged solution retrieval, and the
disciplined mapping of backtracking onto HTTP-level interactions. The
programmer's mental model remains close to that of classical Prolog,
even though transport boundaries intervene.

The \emph{concurrent partition} is the world of actors and interactions.
Computation is structured around long-lived processes equipped with
mailboxes, identities, and lifecycle semantics. Communication is
explicit, asynchronous by default, and potentially distributed across
nodes. Communication-failures and recovery are first-class design concerns.

In the Trinity, the two partitions are deliberately kept orthogonal at the programming model level. A sequential query may be realised as an actor conversation underneath, but the programmer need not know or care; conversely, an actor managing a long-lived protocol is not forced to reason about backtracking and choice points. Each mode has its own transport assumptions, blocking surfaces, and communication-failure semantics.

This separation provides a stable design axis. Systems can be built
primarily in one partition while selectively invoking the other.
Sequential queries may spawn concurrent actors. Actors may internally
perform sequential reasoning. Yet the boundary remains conceptually
intact.

Beyond the partition boundary, several orthogonal axes further structure
the design space. The most important is the axis of \emph{purity}. If
the remote goal in an \texttt{rpc/2-3} call is pure, the call itself is
pure; answers are logical consequences of distributed knowledge bases.
The impure Prolog Web incorporates procedural and extra-logical features
-- cut, dynamic database updates, I/O -- as well as the inherently
procedural character of actor-based concurrency. A related axis is
\emph{openness versus risk}: the extent to which a node exposes itself
as a shared execution environment, with the attendant security and
resource-governance pressures. The profile hierarchy and the capability
posture developed in earlier chapters are best read as disciplined
responses to both axes.

In practice, systems mix pure and impure freely: pure queries for declarative reasoning, actors for coordination and lifecycle management. The guiding principle is \emph{incremental relaxation}:
begin with the purest model that fits the problem, relax purity only
when necessary, and prefer explicit structure over implicit behaviour.
Section~\ref{sec:paradigm-logical-foundation} develops the purity axis
in full.


\subsection{Two levels of abstraction}\label{sec:paradigm-two-levels}

The sequential partition gives rise to a striking property: the Prolog
Web can be read simultaneously at two levels of abstraction. On the
upper level lies a \emph{level of logic programs} -- interlinked Prolog
clauses distributed across nodes, where URIs serve as references between
knowledge bases and \texttt{rpc/2-3} calls serve as the operational
links. On the lower level lies a \emph{level of communicating processes}
-- often realised as toplevel actors that are created, exchange messages,
and eventually terminate.

\begin{figure}[h]
    \centering
	\includegraphics[width=\textwidth]{declarative-procedural}
    \caption{Two levels of abstraction in a homogeneous multi-agent
    system where each agent has access to different knowledge bases.}
    \label{fig:paradigm-decl-proc}
\end{figure}

\noindent On the lower level, processes are created and destroyed and while they
are alive they send messages to each other. Nothing of this kind takes
place on the level of logic programs. Entities such as processes,
messages, and mailboxes do not exist there -- they have been abstracted
away. Drawing a line between these two levels is another way to express
the distinction between declarative and procedural that so often appears
in the story of Prolog. ``What?'' is specified at the level of logic and
the corresponding ``How?'' is handled on the level of communicating
toplevel processes.

The reason this works is that it is very similar to how pure Prolog
explores each branch of a search tree independently from other branches.
Web Prolog, running pure Prolog on the Prolog Web, explores each
\emph{node} independently from other nodes. From the caller's point of
view, a call to \texttt{rpc/2-3} in the body of a rule is just another
subgoal to process, another search tree branch to traverse, only this
time in a Prolog context different from the context of the caller.


This two-level reading is a paradigmatic property. When the sequential partition is pure, the Prolog Web admits both a \emph{logical} reading -- in which distributed computation is understood as inference over distributed knowledge bases -- and a \emph{procedural} reading -- in which the same computation is understood as coordinated message exchange among network processes. Sometimes, when actors and side effects are involved in a fundamental way, only the procedural reading remains. That both readings are available at all in a distributed setting, and that the programmer can move between them, is one of the defining characteristics of the Trinity as a programming paradigm.


\subsection{Unification as the universal data contract}\label{sec:paradigm-unification}

The two-level structure would remain an architectural curiosity if the interface between client and node required the kind of ``contract work'' that characterises conventional web integration. In a typical web API, one reads documentation to learn an endpoint's URL and parameters, inspects a response schema (often JSON with nested fields), writes deserialisation code that maps responses into application structures, and then maintains that glue as the API evolves. The fragility is structural: three logically coupled things are kept separate -- the request shape, the response shape, and the rule that determines whether a response satisfies a request. Conventions such as OpenAPI, JSON Schema, and versioning try to restore this coupling, but only indirectly and after the fact.

In Prolog, unification makes the coupling direct. A goal sent to a node
is simultaneously a request for information, a specification of the
expected structure of answers, and the compatibility test between the
two. The query

\begin{verbatim}
?- employee(Name, Dept, salary(Min, Max)).
\end{verbatim}

\noindent both asks for employees and states that answers must have
exactly this structure. If the node contains
\texttt{employee(alice,\,engineering,\,salary(80000,\,120000))}, the
goal unifies and bindings flow back. If the node only defines \texttt{employee/2}, the query fails immediately: the mismatch is not a downstream runtime error buried in deserialisation code, it is caught at the point of invocation. The query \emph{is} the contract.

Projection is trivial: a client can ask only for what it needs. Evolution is structural: if a node moves from \texttt{employee/2} to \texttt{employee/3}, old clients do not silently receive nulls; their queries simply fail.

Composition is native: constraints combine by conjunction, for example

\begin{verbatim}
?-employee(Name, engineering, salary(_, Max)), Max > 100000
\end{verbatim}

\noindent The combined query can be executed entirely on the node or split between remote and local evaluation. And because a Prolog predicate is inherently relational,
a single relation supports many access patterns depending on which
arguments are instantiated. The client chooses the mode at query time;
the node simply unifies and enumerates. A conventional HTTP API needs
separate endpoints for each of these patterns; in Prolog they are all
the same contract, expressed once, as a goal.


\subsection{Backtracking as pagination}\label{sec:backtracking-pagination}

The structural correspondence between Prolog backtracking and HTTP
pagination is not a surface-level adaptation but a genuine affinity.
Both instantiate the same abstract idea: incremental enumeration under
external control. A paginated endpoint exposes a procedure for generating results and a protocol for consuming them in slices; a Prolog toplevel actor does the same, with choice points serving as the resume mechanism and a ``next'' message as the page-turn signal.

Chapter~\ref{sec:the-prolog-web} developed the operational details: the spectrum from naive re-execution through cached-process resumption to explicit pid-based continuation, and the role of the logical update view in providing snapshot-consistent enumeration. What matters here is the paradigmatic observation. Unification provides the data contract, backtracking provides the enumeration model, and HTTP or WebSockets provides the transport -- with each layer reinforcing the others rather than working against them. This three-way alignment is one of the reasons we regard the Trinity not merely as a practical engineering effort, but as a paradigmatic fit between logic programming and the architecture of the Web.



%%% ───────────────────────────────────────────────────────────────
\section{From compatibility to semantic portability}\label{sec:semantic-portability}

The Prolog community has long understood portability as a matter of \emph{compatibility}: the extent to which different systems happen to accept the same programs and produce similar results. Efforts such as the ISO Prolog standard have addressed this at the level of the core language. Yet, as surveyed in \emph{Fifty Years of Prolog and Beyond}, portability issues remain widespread beyond the core -- in libraries, constraint systems, concurrency models, and interfaces to the outside world. Compatibility is an empirical property; it depends on the contingent alignment of implementations, and it is easily broken by innovation. The Prolog Web sharpens the problem further: in a network of nodes implemented on top of different Prolog systems, the question is no longer just whether two systems accept the same program, but whether they preserve the same observable interactions -- the same patterns of queries, answers, and message exchanges.

In this work, we adopt a different perspective. Rather than treating
portability as compatibility between systems, we treat it as a
\emph{semantic property} of programs.

\begin{quote}
\normalsize \textbf{Semantic portability.} A program is portable across a class of
execution environments if its meaning -- understood as the set of
observable interactions it can produce -- is invariant under relocation
between those environments.
\end{quote}

\noindent This shift has several important consequences. First,
portability becomes relative to a well-defined class of environments: a
design decision, not an accident of history. Second, behaviour rather
than syntax becomes central: two systems may accept the same program and
yet differ in observable behaviour, and under a semantic notion of
portability such differences determine whether portability holds. Third,
portability must be designed rather than inferred: it requires explicit
control over the aspects of execution that affect observable behaviour.

The Trinity instantiates this stronger notion directly. Web Prolog
enforces inexpressibility: the language is restricted so that
client-supplied code cannot depend on uncontrolled or
implementation-specific features. Agents and nodes define observable
interaction: execution is structured around message passing and explicit
interaction, making behaviour visible and comparable across environments.
And node profiles define classes of environments: they provide the
reference points relative to which portability is defined.


\subsection{Node profiles as semantic contracts}\label{sec:node-profiles}

If portability is meaningful only relative to a well-defined class of
execution environments, then those classes must be made explicit. In the
Prolog Trinity ecosystem, this is captured by the notion of a
\emph{node profile}.

\begin{quote}
\normalsize \textbf{Node profile.} A node profile defines a class of execution
environments by specifying the capabilities available to programs and
the semantics of their observable interactions.
\end{quote}

\noindent A node profile is more than a collection of features. It is a
\emph{semantic contract}: it determines what programs may assume about
their execution environment, and what guarantees the environment provides
in return. Two nodes that implement the same profile need not share an
implementation, but they must preserve the same externally visible
behaviour for programs written against that profile.

Profiles play three roles in the Trinity. They define portability
domains: a program is portable precisely within the class of nodes that
implement the profile against which it is written. They separate
semantics from implementation: implementations are free to differ
internally as long as they preserve observable interaction semantics.
And they constrain expressiveness: by limiting available capabilities, a
profile prevents dependence on features that would break portability.

As we saw in Chapter~\ref{sec:prolog-agents}, the Trinity distinguishes
several node profiles, each defining a class of environments over which
portability can be claimed. What varies across profiles is not merely
what a node can do, but what observable behaviour must be preserved when
a program moves between nodes that implement the same profile:

\begin{itemize}
\item \textbf{RELATION.} Observable behaviour consists of the set of
answer substitutions returned for a given query. Two RELATION nodes are
interchangeable if they produce the same answers for the same goals.
No assumptions about execution order, session state, or continuation
need be preserved.

\item \textbf{ISOBASE.} Observable behaviour extends to the
\emph{ordering} and \emph{incremental delivery} of answers. Two ISOBASE
nodes must agree not only on which answers exist, but on the order in
which they are produced and the ability to suspend and resume
enumeration. Portability at this level requires that choice point
semantics and continuation handling are compatible across
implementations.

\item \textbf{ISOTOPE.} Observable behaviour extends to side effects
within a session: I/O, dynamic database updates, and the execution of
client-supplied code. Two ISOTOPE nodes must agree on the available
built-in predicates and on the behaviour of stateful operations --
not merely on what answers are produced, but on what effects are
produced along the way. Portability at this level is the most demanding
within the sequential partition.

\item \textbf{ACTOR.} Observable behaviour is defined in terms of
message exchanges, mailbox semantics, and process lifecycle. Two ACTOR
nodes must preserve the same ordering guarantees on message delivery,
the same failure-propagation semantics, and the same lifetime
discipline for spawned processes. Portability here requires agreement
not about answers to queries, but about the dynamics of concurrent
interaction.
\end{itemize}

\noindent These profiles form a hierarchy of increasingly demanding
portability contracts. At each level, the class of observable behaviour
that must be preserved grows wider, and the semantic agreement required
between implementations grows tighter. What is sufficient for RELATION
-- agreement on answer sets -- is not sufficient for ACTOR, which
requires agreement on the dynamics of interaction itself.

Under this view, portability is no longer a fragile byproduct of partial
standardisation. It becomes a first-class design principle: programs are
written against a profile rather than a specific implementation, their
meaning is defined in terms of observable interactions, and systems are
free to innovate internally as long as they preserve the externally
visible semantics required by the profile.

As shown in Section~\ref{sec:design-space}, when a program is self-contained and the target environment implements the required profile, this principle has a concrete operational payoff: relocation becomes immediate, with no adaptation step required.


%%% ───────────────────────────────────────────────────────────────
\section{Agents as the primary abstraction}\label{sec:paradigm-agents}

At the centre of this design space stands the agent. In the Trinity,
the agent is not an optional library construct layered atop another
abstraction. It is the primary unit of structure.

Actors are agents. Toplevels are agents. Nodes are agents. Statechart
processes are agents in structured disguise. Even external systems --
whether web browsers, services, or AI components -- are treated as
agents when interfacing with the Prolog Web. A node is not merely a
web server that happens to execute Prolog goals; it is a situated
participant that exposes capabilities under a policy, hosts durable
knowledge and services, and mediates between clients, resident actors,
and the surrounding network.

This agent-centric stance contrasts with other dominant structuring
principles. In object-oriented systems, the object is primary. In
functional architectures, the function dominates. In microservice
ecosystems, deployment units often eclipse semantic structure. The
Trinity instead adopts a process-and-message discipline in which
identity, mailbox semantics, and lifecycle policies define the core.
Logical reasoning occurs \emph{within} agents; distribution occurs
\emph{between} agents; behavioural templates constrain \emph{classes} of
agents; and statechart hierarchies govern their reactive control.

The choice has far-reaching implications. It aligns naturally with
distribution-first semantics: because the agent is already a process
with a mailbox, the step from local to distributed execution does not
require a change of abstraction. It supports fault containment and
supervision: agents fail individually, and their failure is observable
through links and monitors. It allows sequential reasoning and reactive
control to coexist: the same agent may answer queries in its sequential
capacity and participate in actor protocols in its concurrent capacity.

The agent-centric stance also supports a literal multi-agent reading.
Once nodes host durable processes that interact over published protocols,
the Prolog Web becomes a web-native multi-agent system in the
architectural sense: a population of autonomous processes coordinating by
message exchange. This is not a claim about rationality or intention in
the AI sense. It is a concrete decomposition principle: each agent has a
mailbox, a lifecycle, and a locally scoped program and database; global
behaviour emerges from explicit communication under failure and timing
constraints.

As Joe Armstrong described the Erlang
approach~\cite{DBLP:conf/hopl/Armstrong07}:

\begin{quote}
Our applications are structured using large numbers of communicating
parallel processes. We take this approach because [\ldots] it provides
an architectural infrastructure -- we can organize our system as a set
of communicating processes. By enumerating all the processes in our
system, and defining the message passing channels between the processes
we can conveniently partition the system into a number of well-defined
sub-components which can be independently implemented, and tested.
\end{quote}

\noindent Armstrong's description applies directly to the Prolog Web. But in the Trinity, the process-and-message discipline is not the whole story -- it is held in place by a tighter structural property.

\subsection{Mutual constraint}\label{sec:paradigm-mutual-constraint}

The three components of the Trinity are mutually constraining -- a property more characteristic of a paradigm than a toolkit. The language profile constrains the agent model: because clients execute
a restricted fragment, agents need not defend against arbitrary
meta-calls, and the security posture can be structural rather than
procedural. The agent model constrains the network substrate: because
agents have explicit identities, mailboxes, and lifecycles, the
interaction protocol must support naming, session tracking, and failure
notification -- requirements that are reflected in the design of the
HTTP and WebSocket layers. And the network substrate constrains the
language profile: because programs must cross transport boundaries, the
language is shaped so that goals, answers, and code fragments are
serialisable as Prolog terms without loss of meaning.

No single component can be redesigned independently without affecting
the others. This tight coupling is not a fragility; it is the source of
the design's coherence. It is also why the Trinity is not well described
as ``Prolog plus actors plus HTTP''. The components do not merely
coexist; they shape each other.

\subsection{Capability-oriented security}\label{sec:capability-security}

One consequence of this mutual shaping is the security model. Because
the language profile, the agent model, and the network substrate
constrain each other, safety does not have to be imposed from outside
--- it arises from the architecture itself.

The structural mechanisms that make the Prolog Web safe to open ---
inexpressibility, invisibility, and containment --- are developed in
Section~\ref{sec:security-prolog-web}. What matters paradigmatically is
that these mechanisms constitute a \emph{capability-oriented} security
model: authority is conveyed by the possession of references, not
granted by external access-control lists. If a program does not hold a
reference to an actor, service, or node, it cannot reach that entity.
The principle can be stated compactly: \emph{possession of a reference
is the authority to use it}.

What is significant is not that the Prolog Web has a security model, but
that its security model is \emph{structural}. Inexpressibility is
enforced at the language boundary: dangerous operations cannot be named
in the client fragment. Invisibility is enforced by the address space: a
pid that has never been revealed cannot be guessed. Containment is
enforced by lifetime dependency: when a session ends, its actor subtree
disappears. Safety arises from the architecture, not from runtime checks
scattered through the code.

This structural character aligns the Prolog Web with a family of
contemporary execution environments --- browser-based JavaScript,
WebAssembly, WASI --- that share the same posture: code executes inside a
deliberately bounded environment and acquires authority only through
explicitly provided references. Traditional Prolog systems are designed
around maximal expressiveness; Web Prolog inverts the default,
treating expressive power as owner-controlled capability that can be
selectively granted rather than silently assumed.


\subsection{Five complementary views}\label{sec:paradigm-five-views}

The agent-centric design refracts into several complementary
perspectives, each emphasising different linked entities and activities:

\begin{description}
\item[\textbf{The Prolog Web of source code.}] URIs link programs to
programs. Nodes expose source code over HTTP; users and tools can ``view
source'' by following links.

\item[\textbf{The Prolog Web of query solutions.}] URIs link callers to
answer streams. A client asks a node to solve a goal and incrementally
enumerates solutions. Interaction is sequential and pull-based.

\item[\textbf{The Prolog Web of actors.}] Pids link live processes with
mailboxes. Interaction is asynchronous: \texttt{!/2} sends,
\texttt{receive/1-2} receives. Linking and monitoring make lifetimes
explicit.

\item[\textbf{The Prolog Web of behaviours.}] Actors implement
behaviours -- protocols specifying message syntax, semantics, and
timing. The behaviour is the contract; the actor is the running entity.

\item[\textbf{The Prolog Web of nodes.}] Nodes are the ``places'' of the
system: they store and serve code, host actors, enforce policy, and
advertise capabilities. They are points of execution and accountability.
\end{description}

\noindent These are not competing architectures but complementary
perspectives on the same system. A Prolog Web application typically
involves all five simultaneously: source code that is browsable, goals
that are queryable, actors that are addressable, behaviours that are
composable, and nodes that are accountable.


%%% ───────────────────────────────────────────────────────────────
\section{The logical foundation: purity, belief, and scoped inference}\label{sec:paradigm-logical-foundation}

The agent-centric view supplies the structural principle; the logical
foundation supplies the semantic one. The Trinity is not merely a
concurrent system that happens to use Prolog syntax. Logic programming
plays a constitutive role, and the key to understanding that role is the
notion of a \emph{pure Prolog Web}.


\subsection{Pure Web Prolog and the pure Prolog Web}\label{sec:paradigm-pure-web-prolog}

As noted in Chapter~\ref{sec:the-prolog-web}, \texttt{rpc/2-3} retains
the logical purity of the predicates it calls. This allows us to form
the notion of \emph{pure Web Prolog}:

\begin{center}
\emph{Pure Web Prolog = Pure Prolog + rpc/2-3}
\end{center}

\noindent Purity here means purity of answer substitutions as logical
consequences; timeouts and connection failures may truncate the answer
stream but do not corrupt it. The union of all partitions of the Prolog
Web written in pure Web Prolog forms what we think of as the \emph{pure
Prolog Web}. With a large number of nodes running pure Web Prolog
programs interlinked to other pure Web Prolog programs, we may \emph{in
theory} create a Web of Logic on top of the conventional Web -- a layer
that can \emph{in principle} grow as large as we want it to grow while
still adhering to the formal semantics of the pure subset of the Prolog
language. In practice, this may never happen, but the \emph{architecture}
is there.

\begin{figure}[h]
    \centering
	\includegraphics[width=8cm]{globe-wrapped-in-logic}
    \caption{A pure Prolog Web spanning the globe remains aspirational, but the architecture does not preclude it.}
    \label{fig:paradigm-globe}
\end{figure}


\subsection{Epistemic reading and scoped inference}\label{sec:paradigm-epistemic}

If agents are the primary abstraction and the pure Prolog Web gives them
a logical character, then a call such as \texttt{rpc(URI,Q)} admits a
natural epistemic reading: it asks whether the agent at \texttt{URI}
\emph{believes} that \texttt{Q}. A clause such as

\begin{verbatim}
human(X) :- rpc('http://n1.org', human(X)).
\end{verbatim}

\noindent can, in the syntax of predicate logic enhanced with a
higher-order predicate \emph{believes/2}, be formulated as

\begin{center}
$\forall x [\mathit{believes}(n_1, \mathit{human}(x)) \rightarrow
\mathit{human}(x)]$
\end{center}

\noindent It makes perfect sense to say that one agent believes that
\texttt{human(socrates)} is true while another does not:

\begin{verbatim}
?- rpc('http://n2.org', human(socrates)).
true.
?- rpc('http://n1.org', human(socrates)).
false.
\end{verbatim}

\noindent And the following query asks whether two agents \emph{agree}
that \texttt{human(X)}:

\begin{verbatim}
?- rpc('http://n1.org', human(X)),
   rpc('http://n2.org', human(X)).
X = plato ;
X = aristotle.
\end{verbatim}

\noindent The epistemic reading is not the only one available. More
conservatively, the URI can be understood as specifying the \emph{scope}
of inference -- the knowledge base with respect to which a query is to
be evaluated. As argued in~\cite{Kifer2005}, scoped inference is
mandatory for realising Negation As Failure on the Web, because applying
any form of the closed world assumption requires knowing the entire
knowledge base first -- something that is not possible for the Web since
its boundaries cannot be clearly delineated. A third and more practical
view is simply that \texttt{rpc/2-3} invokes a goal in the context of a
remote \emph{module}, with data encapsulation that prevents the mixing
of predicates with the same name and arity. For the pure Prolog Web,
the three views -- epistemic operator, scoped inference, remote module
-- converge.

What about \texttt{rpc/3}? Code injection via \texttt{load\_*} options
does not jeopardise purity provided the final composition of the program
executed is itself pure. The \texttt{limit} option is a \emph{pragma}:
it specifies the granularity of the conversation between client and node
but has no effect on the meaning of the query.


\subsection{The less pure Prolog Web}\label{sec:paradigm-less-pure}

Not all of the Prolog Web is pure. ISOBASE nodes may make use of control constructs such as cut, and ISOTOPE nodes go further with side-effecting predicates, I/O, and dynamic database updates. What statelessness constrains is persistent state modification within a single HTTP request -- a protocol-level restriction, not a language-level one.

More fundamentally, the concurrent Prolog Web -- the world of the ACTOR profile -- is inherently procedural. Actor-based programming involves process creation and destruction, message ordering, selective receive, and failure propagation -- none of which have declarative readings. A realistic application will typically mix profiles freely, using pure queries for declarative reasoning and actors for coordination and state.

This spectrum from pure to impure is itself a paradigmatic feature.
Rather than forcing a choice between declarative and procedural -- a
tension that has shadowed Prolog since its inception -- the Trinity
makes the boundary explicit and navigable. The programmer can reason
about which parts of a system admit logical readings and which require
procedural ones, and can move between the two with architectural
awareness rather than accidental drift.


%%% ───────────────────────────────────────────────────────────────
\section{The distribution-first model}\label{sec:paradigm-distribution-first}

Web Prolog adopts a distribution-first view: concurrency and
distribution are not anomalies to be managed but the fundamental reality
of computation. Local execution is simply the case where network latency
happens to be zero. The Actor Model views computation as inherently
distributed, with independent actors communicating asynchronously;
sequential processing is a special case. For Web Prolog, treating
distribution as normal rather than exceptional simplifies the programming
model and aligns naturally with web-scale deployment.

The formal foundation for this stance comes from
\citet{conf/erlang/SvenssonFE10}, who proposed a unified semantics for
Erlang-like systems. Their key insight is to eliminate the distinction
between local and remote communication: all messages travel through a
single conceptual mechanism, and side-effecting operations are handled
uniformly by a node controller abstraction. This makes placement a
configuration detail rather than a semantic distinction. Web Prolog can
adopt this cleaner semantic story from the outset, rather than being
constrained by decades of backward compatibility as Erlang is.
Svensson, Fredlund, and Benac Earle themselves hinted at this when they
noted that ``perhaps the language for this semantics is not Erlang,
maybe there is another language waiting around the corner.'' We believe
the language waiting around the corner is Web Prolog.

A useful way to read the distribution-first stance is through the
classic goal of \emph{distribution transparency}: making the
distribution of processes and resources invisible to applications. The
Prolog Web does not aim for transparency by hiding the network behind a
magic API; rather, it aims for a disciplined form in which the
\emph{same} conceptual operations --- \texttt{spawn}, send, monitoring,
\texttt{rpc/2-3} --- apply locally and remotely. Latency and
communication-failure probabilities change with distance, but the
programming model does not bifurcate into a ``local'' and a ``remote''
language. What makes this disciplined rather than naive is the profile
hierarchy: the profiles specify exactly which guarantees hold across
the local/remote boundary, so that transparency is a known contract
rather than a leaky abstraction.

The combination of Prolog's execution model with this actor semantics
can be given a clean, modular formal account. Conceptually, we extend
the unified actor semantics with a disjoint state component for each
process: in addition to its mailbox and process control state, an actor
carries a \emph{Prolog engine state} consisting of a goal stack and a
current substitution. The Erlang-style rules manipulate only processes
and mailboxes; the Prolog rules touch only goals and substitutions.
Because these state components are disjoint, the two rule sets can be
interleaved without conflict. Backtracking remains local to a single
actor's engine; interaction between actors is always mediated by
messages. A full development of this layered semantics is beyond the
scope of this chapter, but the composability of the two rule sets is
what justifies treating the sequential and concurrent partitions as
genuinely orthogonal.


\subsection{The browser as a Web Prolog runtime}\label{sec:paradigm-browser}

If placement is a configuration detail, then the distribution-first
model extends naturally to the client side. A browser equipped with a
Web Prolog runtime is not a node -- it does not accept incoming
connections -- but it is a first-class participant in the Prolog Web,
able to use the same interaction primitives available to server-side
code. The programmer should be able to use synchronous calls when
blocking is acceptable, non-blocking request--response calls when
responsiveness matters, and actor-style interaction for cases that
genuinely require long-lived conversational state -- all without being
exposed to transport-level details.

The longer-term ambition suggested by this design is Prolog on both
sides of the Web. At the same time, this should not be confused with an attempt to hide
the Web or to force all clients into a single language. Many developers
will continue to prefer writing front-ends in JavaScript and
communicating with Prolog back-ends through explicit web APIs. The
Prolog Web should make such APIs visible, well-specified, and robust. A
Prolog-everywhere experience may make them largely invisible in the
common case, but engineering robustness requires that they remain
accessible.


%%% ───────────────────────────────────────────────────────────────
\section{Control as a first-class dimension}\label{sec:paradigm-control}

Consider the famous equation coined by Robert Kowalski:

\begin{quotation}
\normalsize\noindent Algorithm $=$ Logic $+$ Control
\end{quotation}

\noindent The kind of control we have in mind here is different from
search control. Conversational and reactive systems are hard not
primarily because they require sophisticated algorithms, but because
they require fine-grained \emph{interaction control}: turn-taking,
timeouts, interruption, recovery, orchestration, and long-lived
dialogue flow. In the Trinity, this motivates the combination of three
programming models: logic programming for knowledge and inference;
actor-based programming for concurrency, distribution, and interaction;
and statecharts for explicit, inspectable control structure. The three
are orthogonal: each addresses a dimension of system design that the
others leave open.


\subsection{Governance of interaction}\label{sec:paradigm-governance}

If the Trinity were only a story about distributed queries and actor
concurrency, it would leave an important practical question unanswered:
how does one \emph{govern} complex interaction so that it remains
inspectable, auditable, and robust? A great deal of complexity in web
applications and agent systems is about \emph{control}: protocol
sequencing, turn-taking, timeouts, retries, cancellation, interruption,
recovery, and the disciplined handling of exceptional cases.

Statechart actors are the Trinity's response. The idea is not to replace
actor programs with statecharts, but to offer a structured substrate for
reactive behaviour in which states, transitions, events, timers, guards,
and recovery paths are explicit and therefore reviewable. Statechart
actors turn the actor model from a concurrency substrate into a
\emph{governable interaction platform}.

On the level of notation, this can be realised in more than one way: a
Prolog-friendly representation for SCXML-style charts, with translation
to and from SCXML when interchange is useful; or SCXML itself as the
control notation with Web Prolog as the language of guards, actions, and
data. Either way, the relationship follows a familiar Web pattern: a
declarative control layer whose behaviour is scripted in a
general-purpose language.


\subsection{A platform for the Agentic Web}\label{sec:paradigm-agentic-web}

Chapter~\ref{sec:prolog-ai-agents-on-the-agentic-web} placed the
Trinity in a broader movement: software is becoming increasingly
agentic. The Trinity's claim is that the Prolog Web is a natural
substrate for such systems because it integrates three capabilities that
agentic applications repeatedly need: a web-native network of nodes, a
concurrency model for durable interactions, and a symbolic language for
explicit representation and constraint.

The Agentic Web framing supplies the \emph{application pressure}:
conversational and tool-using agents need explicit protocol state,
recovery from partial failure, and coordination under bounded waiting.
Statecharts supply the \emph{governance mechanism}: inspectable control
with timeouts, retries, and fallback paths as first-class structure.
Logic programming supplies the \emph{symbolic substrate}: knowledge,
constraints, and goal-directed inference remain available inside agents
and can be shipped, composed, and queried across node boundaries.

This synthesis motivates a pragmatic neuro-symbolic stance. Statistical
language models are useful components in agentic systems, but they are
not reliable as autonomous governors of protocol and state. The Trinity
treats statistical components, when present, as powerful but fallible
modules, and places responsibility for durable behaviour in explicit
symbolic structures. ``Agentic'' does not merely mean ``chatty''. It
means \emph{durably interactive under explicit rules}, in an open
network where failure and delay are normal.



%%% ───────────────────────────────────────────────────────────────
\section{Executable logic on the Web}\label{sec:paradigm-executable-logic}

A recurring posture in web architecture is to treat logic as primarily
\emph{representational}: a way to publish structured meaning, express
constraints, and support inference over shared data. The Semantic Web
strand exemplifies this orientation. It emphasises common identifiers,
interoperable data models, and well-defined entailment regimes. These
are substantial achievements, but they leave open a practical question:
where does executable reasoning live, and how is it composed across
independently hosted sites?

The Trinity's distinctive move is to treat logic on the Web as
\emph{executable} in a disciplined sense. A node is not only a publisher
of facts; it is a host for relational services, rules, and interaction
protocols that can be invoked, combined, and enumerated as computations.
In the sequential partition, a goal can be shipped to a node and
answered nondeterministically, with unification as the data contract and
backtracking as the enumeration model. In the concurrent partition,
long-lived agents coordinate by message exchange while embedding
explicit symbolic structure in the messages they exchange and the
constraints they apply.

Once this substrate exists, ``logic on the Web'' is no longer confined
to representation and offline reasoning. It becomes a first-class
component of live systems: callable, composable, inspectable, and
governable. This executable substrate is what makes it natural to treat
the Trinity as a form of symbolic middleware for the Web.


\subsection{Positioning relative to the Semantic Web}\label{sec:paradigm-positioning-sw}

The Prolog Web and the Semantic Web are similar in spirit: both treat
the Web as a computational environment, both rely on global naming and
logic-based machinery, and both aim at decentralisation. But they
emphasise different things. The Semantic Web places its centre of
gravity on interoperable data models and shared semantic commitments;
the Prolog Web places its centre of gravity on executable relational
services, process-level structure, and web-native agents.

\begin{table}[ht]
\begin{center}
\begin{tabular}{ | p{5.7cm} | p{5.7cm} | }
\hline
\textbf{The Prolog Web} & \textbf{The Semantic Web} \\ \hline
Uses the same language for querying and rule-based computation &
Typically separates data representation (RDF/OWL) from querying
(SPARQL) and application logic \\ \hline
Provides a process and service architecture as part of the proposal &
Provides standards for interoperable data and ontology semantics, but
does not prescribe a process architecture \\ \hline
Supports agent-style computation through long-lived actors and message
passing & Has strong support for shared ontologies and open-world
reasoning through standardised languages \\ \hline
Shared meaning is typically implicit in executable predicates; shared
ontologies are optional & Interoperability is centred on shared
identifiers, vocabularies, and ontology constraints \\ \hline
Well suited as a programming and orchestration layer around
heterogeneous services & Well suited as a lingua franca for knowledge
interchange \\ \hline
\end{tabular}
\end{center}
\caption{Comparing the Prolog Web and the Semantic Web.}
\label{tab:paradigm-wp-sw-comparison}
\end{table}

\noindent The Semantic Web asks: \emph{how can independently published
data be understood and combined?} The Prolog Web asks: \emph{how can
independently hosted logic be executed and composed?} The two questions
are complementary. RDF, OWL, and SPARQL provide machinery for
representing and querying structured knowledge. The Trinity adds an
execution substrate: a way to run queries, host rules, ship code, and
maintain long-lived interaction under explicit operational rules. The
paradigmatic claim is that the Trinity fills a gap the Semantic Web has
long acknowledged: the gap between representation and execution, between
knowing and doing.

Although this section uses the Semantic Web as a concrete reference
point, most of the underlying issues are not specific to RDF, OWL, or
W3C standards. Property-graph systems and large-scale knowledge-graph
platforms have converged on similar concerns. The Prolog Web addresses
this broader design space, offering an executable logic-based substrate
that can sit alongside different graph representations rather than
competing with any single one. This is consistent with the multi-stack
architecture for the Semantic Web proposed
in~\citep{Kifer2005,Horrocks2005}.


%%% ───────────────────────────────────────────────────────────────
\section{Comparison to other approaches}\label{sec:paradigm-related-work}

The claim that the Trinity constitutes a candidate paradigm invites a
natural sceptical response: has this not been attempted before? Three
precedents are close enough to warrant explicit comparison.

\subsection{Linda and tuple-space coordination.}
Linda~\citep{gelernter1985generative} introduced coordination through a
shared associative memory, with pattern matching over tuples as the
retrieval mechanism. The Trinity shares the emphasis on pattern matching
as coordination, but the Prolog Web is a network of active agents with
explicit identities and lifecycles, not a passive shared store.
Coordination is directional: a client submits a goal to a named node
and receives answers through a well-defined protocol. This directedness
makes reasoning about authority, lifetime, and failure tractable.

\subsection{Concurrent Prolog and the committed-choice languages.}
Concurrent Prolog~\citep{shapiro1983subset},
Parlog~\citep{clark1986parlog}, and GHC~\citep{ueda1985guarded} showed
that SLD-resolution and concurrency are not incompatible, but they paid
a significant price: the declarative reading is severely weakened by the
commitment mechanism, and backtracking is absent from the concurrent
partition. The Trinity keeps both. The sequential partition retains full
Prolog semantics; the concurrent partition adopts Erlang-style message
passing, making the non-declarative character of concurrency explicit.
The two are deliberately orthogonal: backtracking does not roll back
mailboxes, and actor communication does not interfere with proof search.

\subsection{Oz/Mozart and the multiparadigm kernel language.}
Oz~\citep{smolka1995oz}, implemented in the Mozart
system~\citep{vanroy2003logic}, is the closest and most instructive
precedent. Like the Trinity, it combines logic programming, concurrency,
and distribution in a single coherent design. Distributed
Oz~\citep{vanroy1997distributed} adds network-transparent mobility and
fault tolerance, and the van Roy--Haridi
textbook~\citep{vanroy2004concepts} provides an unusually clean formal
account through a hierarchy of kernel languages. The comparison
therefore cannot rest on missing features; it must rest on design
stance.

The decisive difference is that Oz replaces Horn clause syntax with a
fully compositional, higher-order kernel language that accommodates all
its paradigms uniformly. The result is a coherent and expressive design,
but one that is no longer Prolog in any meaningful sense: it abandons
unification-driven clause selection, backtracking, and the program-as-database
identity that gives Prolog its distinctive character. The Prolog
community has not adopted Oz, and the distance is not bridgeable by a
compatibility layer. The Trinity, by contrast, is defined as a
\emph{profile} of Prolog. It extends ISO Prolog with concurrency
primitives and distribution infrastructure while remaining syntactically
and semantically recognisable to anyone who already knows the language.
This is a deliberately narrower ambition -- but it is the right ambition
for a proposal that aims to unify an existing fragmented community
rather than replace it with a new one.

A second difference concerns the Web. Oz/Mozart achieves genuine network
transparency through a custom protocol over TCP, but nodes are not URIs,
processes are not addressable by HTTP clients, and services are not
discoverable by web tooling. The Trinity takes the opposite stance: HTTP
and WebSocket are the primary transport substrate, and the
addressability of nodes by URI is a design commitment. The result is a
system less general than Oz in the space of distribution models but more
naturally embedded in the Web as it actually exists.


%%% ───────────────────────────────────────────────────────────────
\section{Why the Trinity fits Prolog}\label{sec:why-prolog-trinity}

The Trinity architecture is not merely implemented \emph{in} Prolog; it
resonates with properties of the language that make the fit unusually
tight.

Unification gives the ecosystem a universal data contract: structured
terms serve simultaneously as data structures, messages, and query
templates, and the same mechanism that matches a goal against a clause
head also matches incoming messages against expected patterns.
Nondeterministic computation provides a disciplined enumeration model
that maps naturally onto Web-style incremental consumption. The
proximity between code and data means that predicate shipping and remote
goal execution can be expressed using ordinary language constructs rather
than elaborate serialisation frameworks. And Prolog's strengths in
symbolic reasoning and explanation-oriented computation complement the
broader ambition of supporting agents capable of reasoning about
structured knowledge.

What tends to disappear when the pattern is transplanted to other
languages is the degree of \emph{semantic unity}. In Prolog, terms,
clauses, goals, and substitutions form a single conceptual framework
supporting knowledge representation, communication, and computation
simultaneously. In most other languages these concerns are handled by
separate mechanisms: data structures, serialisation formats, rule
engines, and orchestration frameworks. The resulting systems may be
powerful, but they lack the conceptual compression that Prolog provides.
The Trinity is therefore best viewed as a web-native expression of the
logic programming paradigm -- the paradigm in which Prolog has
historically found its most direct realisation.


%%% ───────────────────────────────────────────────────────────────
\section{Summary}\label{sec:paradigm-summary}

This chapter has stepped back from the technical machinery developed in
earlier chapters to ask what conceptual pattern emerges when the
components of the Prolog Trinity are considered as a whole.

The answer is a candidate paradigm for web-native symbolic computation,
characterised by several interlocking and mutually constraining
commitments. Two orthogonal partitions -- sequential and concurrent --
provide a stable design axis that separates declarative reasoning from
actor-based interaction. Semantic portability and node profiles replace
fragile compatibility with a designed, contract-based notion of what it
means for programs to retain their meaning across environments. The
agent serves as the primary abstraction, unifying actors, toplevels,
nodes, and statechart processes under a single process-and-message
discipline -- and the five complementary views of the Prolog Web show
how this abstraction refracts into a rich but coherent set of
perspectives. A logical foundation preserves the possibility of pure,
declarative reasoning even in a distributed setting, with the epistemic
reading of \texttt{rpc/2-3} and scoped inference giving formal substance
to what it means to query a web of knowledge bases. Control, in the
sense of interaction governance, is treated as a first-class design
dimension through the integration of statecharts alongside logic and
actors. And the mutual constraint among language profile, agent model,
and network substrate -- each shaping and being shaped by the others --
is what gives the design its coherence and prevents it from decomposing
into a loose collection of independent ideas.

Whether this constellation of commitments constitutes a paradigm in the
strongest sense of the word is ultimately a judgment that the community,
not the author, must make. What we can say is that these commitments are
not accidental. They emerged from a deliberate attempt to integrate
logic programming, actor-based concurrency, and web architecture into a
coherent whole -- and that the result, for better or worse, does not
reduce to any of its components.